\documentclass[11pt,a4paper]{article}
\usepackage[utf8]{inputenc}
\usepackage[spanish]{babel}
\usepackage{geometry}
\geometry{top=2.5cm, bottom=2.5cm, left=2.5cm, right=2.5cm}
\usepackage{hyperref}
\usepackage{graphicx}
\usepackage{listings}
\usepackage{xcolor}
\usepackage{tikz}
\usetikzlibrary{shapes.geometric, arrows, positioning}

\definecolor{codegreen}{rgb}{0,0.6,0}
\definecolor{codegray}{rgb}{0.5,0.5,0.5}
\definecolor{codepurple}{rgb}{0.58,0,0.82}
\definecolor{backcolour}{rgb}{0.95,0.95,0.92}

\lstdefinestyle{mystyle}{
    backgroundcolor=\color{backcolour},   
    commentstyle=\color{codegreen},
    keywordstyle=\color{magenta},
    numberstyle=\tiny\color{codegray},
    stringstyle=\color{codepurple},
    basicstyle=\ttfamily\footnotesize,
    breakatwhitespace=false,         
    breaklines=true,                 
    captionpos=b,                    
    keepspaces=true,                 
    numbers=left,                    
    numbersep=5pt,                  
    showspaces=false,                
    showstringspaces=false,
    showtabs=false,                  
    tabsize=2
}
\lstset{style=mystyle}

\title{\textbf{Informe Técnico del Compilador HULK}}
\author{Michell Viu Ramirez}
\date{\today}

\begin{document}

\maketitle
\tableofcontents
\newpage

\section{Introducción}
HULK (Havana University Language for Kompilers) es un lenguaje de programación funcional orientado a objetos diseñado con fines educativos. El compilador actual, escrito completamente en \textbf{Rust}, consta de un pipeline robusto que procesa el código fuente de HULK, genera un Árbol de Sintaxis Abstracta (AST), realiza un exhaustivo chequeo semántico y compila hacia código nativo utilizando \textbf{LLVM} a través de la interfaz de la librería \texttt{inkwell}.

\section{Estructura del Proyecto}
El entorno de desarrollo está organizado para modularizar el proceso de compilación separando la herramienta de parsing, la lógica del compilador y el entorno de ejecución (runtime):

\begin{itemize}
    \item \texttt{src/parser/}: Almacena toda la lógica de reconocimiento sintáctico, empleando \texttt{LALRPOP}. Contiene definiciones del AST (\texttt{ast/}), el chequeo semántico (\texttt{semantic/}), y los tokens sintácticos (\texttt{tokens/}).
    \item \texttt{src/codegen/}: Se encarga de traducir el AST verificado al Intermediate Representation (IR) de LLVM. Contiene submódulos como \texttt{builtins}, \texttt{classes}, \texttt{context}, \texttt{expressions}, etc.
    \item \texttt{runtime/}: Código C suplementario (\texttt{hulk\_runtime.c}) compilado estáticamente junto con el IR de HULK. Provee el soporte mínimo en tiempo de ejecución, por ejemplo funciones de E/S, gestión de tipos y reservado de memoria.
    \item \texttt{tests/}: Directorio de scripts \texttt{.hulk} que valida el correcto funcionamiento de las pasadas del compilador.
\end{itemize}

\section{Gramática y Análisis Sintáctico}

El compilador utiliza un generador de parsers \textbf{LR(1)/LALR(1)} llamado \textbf{LALRPOP}. La gramática implementa directamente la explícita jerarquía de precedencias evitando conflictos shift/reduce, asegurando una semántica que no requiere ambigüedad por "dangling-else" ni problemas de asociatividad.

\subsection{Tipos y Nodos}
Todos los constructos son expresiones, garantizando un flujo orientado al comportamiento funcional. El análisis incluye la definición local (\texttt{let}), condicionales (\texttt{if}), iteración (\texttt{while}, \texttt{for}), chequeos de tipo en tiempo de ejecución (\texttt{is}) e instanciaciones de clases tipadas dinámicamente o arreglos (\texttt{new}). Recientemente se adoptaron los identificadores modernos: \texttt{type} para inicializar clases y \texttt{inherits} para las declaraciones de herencia.

\section{Flujo Interno del Sistema}

A continuación, se ilustra el flujo base de una compilación estándar:

\begin{center}
\begin{tikzpicture}[
    node distance=2cm and 1cm,
    block/.style ={rectangle, draw, fill=blue!10, text width=9em, text centered, rounded corners, minimum height=3em},
    line/.style ={draw, -latex'}
]

\node [block] (source) {Código HULK (\texttt{.hulk})};
\node [block, below of=source] (lexer) {LALRPOP (Lexer)};
\node [block, right of=lexer, node distance=5cm] (parser) {LALRPOP (Parser)};
\node [block, below of=parser] (ast) {AST en Rust};
\node [block, left of=ast, node distance=5cm] (semantic) {Chequeo Semántico \& Inferencia de Tipos};
\node [block, below of=semantic] (codegen) {Generación de Código (LLVM IR)};
\node [block, right of=codegen, node distance=5cm] (runtime) {C Runtime (Clang/GCC)};
\node [block, below of=codegen] (linker) {Linker (Binario final)};

\path [line] (source) -- (lexer);
\path [line] (lexer) -- (parser);
\path [line] (parser) -- (ast);
\path [line] (ast) -- (semantic);
\path [line] (semantic) -- (codegen);
\path [line] (codegen) -- (linker);
\path [line] (runtime) |- (linker);

\end{tikzpicture}
\end{center}

\newpage
\section{Descripción Detallada de Etapas}

\subsection{Lexer y Parser}
La fase de Lexer ocurre automáticamente gracias a las expresiones regulares embebidas directamente en el archivo \texttt{grammar.lalrpop}. Reconoce tokens de \texttt{NUMBER}, \texttt{STRING}, palabras reservadas (\texttt{type}, \texttt{if}, \texttt{for}, etc.) e ignora interlineados. Luego, el parser consume estos tokens organizados estrictamente por niveles. LALRPOP reporta un error estructurado si el código ingresado difiere de alguna regla de precedencia en la jerarquía (ej. de precedencia mayor: llamadas de método y accesos de arrays, precedencia menor: declarativas de \texttt{let}).

\subsection{Chequeo Semántico}
\begin{itemize}
    \item \textbf{Collector}: Realiza una primera pasada superficial reconociendo qué tipos/clases han sido declaradas globalmente para alimentar una Tabla de Símbolos en crudo.
    \item \textbf{Builder}: Examina la estructura de clases (padres e hijos, \texttt{inherits}) y captura las definiciones globales de las firmas de métodos y funciones (parámetros y tipo de retorno).
    \item \textbf{Chequeo final}: Recorre el AST validando ámbitos (scoping), uso lícito de variables e integridad estructural. 
\end{itemize}

\subsection{Inferencia de Tipos y Chequeo Fuerte}
Puesto que HULK es fuertemente tipado estáticamente, cada expresión expuesta obtiene un retorno derivado. 
El sistema resuelve progresivamente tipos implícitos. Si un bloque \texttt{if} tiene múltiples ramas, infiere el Ancestro Común Más Cercano (LCA) entre los tipos retornados por cada bloque. Se asegura localmente de validar que el tipo derecho al momento de una asignación destructiva (\texttt{:=}) sea compatible o igual al tipo destino. Finalmente la existencia del operador \texttt{is} permite saltar estas verificaciones validando referencias dinámicamente.

\subsection{Generación de Código Base (Codegen)}
El módulo \texttt{codegen} convierte el AST decorado, ahora validado, hacia instrucciones de LLVM. Por medio del Wrapper \texttt{inkwell}:
\begin{itemize}
    \item Se compilan modelos de \texttt{VTable} para simular el \textit{dispatch} dinámico de clases (Polimorfismo).
    \item Las estructuras de control generan bloques básicos en LLVM IR enlazados por \texttt{br} (branch instructions). 
    \item Operaciones \texttt{new} y creación de \textit{Strings} derivan la lógica hacia funciones en el runtime de C como \texttt{malloc}.
\end{itemize}

\subsection{Entorno de Ejecución (Runtime)}
Dado que el binario compila de forma nativa al microprocesador host, las rutinas base están empaquetadas en un módulo en \texttt{C} en \texttt{runtime/hulk\_runtime.c}. En su interior implementa el Print de la consola fundamental de HULK, concadenado de texto Unicode a primitivas simples de `char*` en C plano y rutinas auxiliares necesarias por el recolector de memoria en fases dinámicas del ciclo de ejecución. 

\section{Conclusiones}
La refactorización con un stack moderno de LALRPOP simplifica en gran medida el mantenimiento a futuro. Las implementaciones LLVM garantizan que independientemente de la carga funcional del lenguaje, los binarios emitidos resultarán altamente optimizados y libres del peso computacional clásico de una Máquina Virtual.

\end{document}